%% ma: L12-B18-0006
%% lop: 12
%% chuong: 6
%% bai: 18
%% dang: 12.18.5
%% loai: TLN
%% muc: VD
%% dap_an: "0,49"
%% nguon: "(NBV)-ĐỀ SỐ 13-MỖI NGÀY 1 ĐỀ THI 2025.pdf (D13, MỖI NGÀY 1 ĐỀ - Đề số 13), Phần III Câu 4"
%% kien_thuc: "Bài 18 (xác suất có điều kiện), công thức cộng xác suất và nhân xác suất của hai biến cố độc lập (lớp 11)"
%% trang_thai: cho-duyet
%% ly_do: "Dạng mới 12.18.5 (xác suất một xạ thủ trúng khi biết có ít nhất một viên trúng đích) chờ giáo viên duyệt"
%% ghi_chu: "Đáp án gốc 0,49 đúng (P(Y trúng | có viên trúng đích) = 0,4/0,82 = 0,4878). Đã bỏ ảnh minh họa xạ thủ."
\begin{ex}
Hai xạ thủ $X$ và $Y$ mỗi người bắn một viên đạn vào cùng mục tiêu độc lập. Giả sử xác suất bắn trúng đích của hai xạ thủ $X$ và $Y$ lần lượt là $0{,}7$ và $0{,}4$. Biết rằng có viên đạn trúng đích, tính xác suất để $Y$ bắn trúng đích (kết quả làm tròn đến hàng phần trăm).
\shortans{0,49}
\loigiai{Gọi $A$: ``$X$ bắn trúng'', $B$: ``$Y$ bắn trúng''; $A$, $B$ độc lập, $P(A)=0{,}7$, $P(B)=0{,}4$.

Biến cố ``có viên đạn trúng đích'' là $A\cup B$ với $P(A\cup B)=P(A)+P(B)-P(A)P(B)=0{,}7+0{,}4-0{,}28=0{,}82$.

Vì $B\subset A\cup B$ nên $P\left(B\cap(A\cup B)\right)=P(B)=0{,}4$. Vậy $P\left(B\mid A\cup B\right)=\dfrac{0{,}4}{0{,}82}\approx0{,}49$.}
\end{ex}
